% Part: first-order-logic
% Chapter: sequent-calculus
% Section: quantifier-rules

\documentclass[../../../include/open-logic-section]{subfiles}

\begin{document}

\olfileid{fol}{seq}{qrl}

\olsection{संख्यापकांचे नियम}

\subsection{$\lforall$ साठीचे नियम}

\begin{defish}
\Axiom$ !A(t), \Gamma \fCenter \Delta$
\RightLabel{\LeftR{\lforall}}
\UnaryInf$ \lforall[x][!A(x)],\Gamma \fCenter \Delta$
\DisplayProof
\hfill
\Axiom$ \Gamma \fCenter \Delta, !A(a) $
\RightLabel{\RightR{\lforall}}
\UnaryInf$ \Gamma \fCenter \Delta, \lforall[x][!A(x)]$
\DisplayProof
\end{defish}

\LeftR{\lforall} मध्ये, $t$ हे बंद पद आहे (म्हणजे, चरे नसलेले पद).
\RightR{\lforall} मध्ये, $a$~हा असा !!a{constant} आहे की तो
\RightR{\lforall} नियमाच्या खालच्या क्रमवर्तीमध्ये कोठेही येता कामा नये.
\RightR{\forall} अनुमानाचा $a$~हा \emph{आयगेन चर} आहे, असे आपण
म्हणतो.\footnote{वरील नियमात $a$ हा !!a{constant} असला तरी आपण
``आयगेन चर'' ही संज्ञा वापरतो. यामागे ऐतिहासिक कारणे आहेत.}

\subsection{$\lexists$ साठीचे नियम}

\begin{defish}
\Axiom$ !A(a), \Gamma \fCenter \Delta $
\RightLabel{\LeftR{\lexists}}
\UnaryInf$ \lexists[x][!A(x)], \Gamma \fCenter \Delta$
\DisplayProof
\hfill
\Axiom$ \Gamma \fCenter \Delta, !A(t) $
\RightLabel{\RightR{\lexists}}
\UnaryInf$ \Gamma \fCenter \Delta, \lexists[x][!A(x)]$
\DisplayProof
\end{defish}

पुन्हा, $t$~हे बंद पद आहे आणि $a$~हा असा !!a{constant} आहे की तो
\LeftR{\lexists} नियमाच्या खालच्या क्रमवर्तीमध्ये येत नाही.
\LeftR{\lexists} अनुमानाचा $a$~हा \emph{आयगेन चर} आहे, असे आपण म्हणतो.

\RightR{\lforall} किंवा \LeftR{\lexists} अनुमानाच्या खालच्या
क्रमवर्तीमध्ये आयगेन चर येत नसावा, या अटीला \emph{आयगेन चराची अट}
म्हणतात.

\begin{explain}
ही प्रथा आठवा: $!A$ हे असे !!a{formula} असेल की त्यात
!!{variable}~$x$ मुक्त असेल, तर हे आपण~$!A(x)$ असे लिहून दर्शवतो.
त्याच संदर्भात, मग $!A(t)$ हा~$\Subst{!A}{t}{x}$ चा संक्षेप असतो.
त्यामुळे \RightR{\lexists} नियम पुढीलप्रमाणेही लिहिता येईल:
\begin{prooftree}
  \Axiom$\Gamma \fCenter \Delta, \Subst{!A}{t}{x}$
  \RightLabel{\RightR{\lexists}}
  \UnaryInf$\Gamma \fCenter \Delta, \lexists[x][!A]$
\end{prooftree}
लक्षात घ्या की~$!A$ मध्ये $t$ आधीच येऊ शकते; उदा., $!A$ हे
~$\Atom{\Obj P}{t,x}$ असू शकते. म्हणून,~$\Gamma \Sequent \Delta,
\Atom{\Obj P}{t,t}$ पासून $\Gamma \Sequent \Delta,
\lexists[x][\Atom{\Obj P}{t,x}]$ निष्पन्न करणे हा
\RightR{\lexists} चा योग्य उपयोग आहे---आधारविधानात~$t$ ज्या ठिकाणी
आले आहे, त्यांपैकी एक किंवा अधिक ठिकाणी (म्हणजे, त्या पदाच्या एक किंवा अधिक
आस्थितींच्या जागी) बद्ध !!{variable}~$x$ घालता येतो; सर्वच आस्थिती बदलणे आवश्यक नाही. परंतु
\RightR{\lforall} आणि~\LeftR{\lexists} मधील आयगेन चराच्या अटींनुसार
!!{constant}~$a$ हा~$!A$ मध्ये येता कामा नये. म्हणून,
$\RightR{\lforall}$ वापरून $\Gamma \Sequent \Delta, \Atom{\Obj P}{a,a}$
पासून $\Gamma \Sequent \Delta, \lforall[x][\Atom{\Obj P}{a,x}]$
योग्य रीतीने निष्पन्न करता येत नाही.
\end{explain}

\begin{explain}
\RightR{\lexists} आणि \LeftR{\lforall} मध्ये पद~$t$ बंद असण्याव्यतिरिक्त
त्यावर आणखी कोणतेही निर्बंध नाहीत. याउलट, \LeftR{\lexists} आणि \RightR{\lforall}
नियमांमध्ये आयगेन चराच्या अटीनुसार वरच्या क्रमवर्तीमध्ये~$!A(a)$ च्या
बाहेर !!{constant}~$a$ कोठेही येता कामा नये. ही पद्धत निर्दोष आहे,
म्हणजे ती केवळ वैध क्रमवर्तीच !!{derive}s करते, याची खात्री करण्यासाठी ही
अट आवश्यक आहे. ही अट नसती, तर पुढील निष्पत्त्या अनुज्ञात असत्या:
\begin{prooftree}
  \Axiom$!A(a) \fCenter !A(a)$
  \RightLabel{*\LeftR{\lexists}}
  \UnaryInf$\lexists[x][!A(x)] \fCenter !A(a)$
  \RightLabel{\RightR{\lforall}}
  \UnaryInf$\lexists[x][!A(x)] \fCenter \lforall[x][!A(x)]$
  \DisplayProof\bottomAlignProof
  \qquad
  \Axiom$!A(a) \fCenter !A(a)$
  \RightLabel{*\RightR{\lforall}}
  \UnaryInf$!A(a) \fCenter \lforall[x][!A(x)]$
  \RightLabel{\LeftR{\lexists}}
  \UnaryInf$\lexists[x][!A(x)] \fCenter \lforall[x][!A(x)]$
\end{prooftree}
परंतु $\lexists[x][!A(x)] \Sequent \lforall[x][!A(x)]$ ही क्रमवर्ती
वैध नाही.
\end{explain}

\end{document}
